\begin{proofbox}

	\textbf{\textcolor{CProof}{思路提示}}

	取椭圆短轴端点，利用椭圆定义到两焦点距离之和为 $2a$，通过距离公式计算即可得到 $a^{2}=b^{2}+c^{2}$。

	\medskip
	\textbf{\textcolor{CProof}{正式推导}}
	\begin{description}[style=nextline, leftmargin=2.8em, labelsep=0.5em, itemsep=0.55em]
		\item[\textbf{步骤一（设椭圆方程）：}]
			设椭圆标准方程为 $\frac{x^{2}}{a^{2}}+\frac{y^{2}}{b^{2}}=1$（$a>b>0$），焦点为 $F_{1}(-c,0)$、$F_{2}(c,0)$。
			\[
				\frac{x^{2}}{a^{2}}+\frac{y^{2}}{b^{2}}=1,
				\qquad
				F_{1}(-c,0),\;F_{2}(c,0).
			\]
			\par\textbf{理由：}
			这是椭圆标准形式，焦点在 $x$ 轴上（焦点在 $y$ 轴时推导类似）。

		\item[\textbf{步骤二（取特殊点）：}]
			取椭圆上一点 $P(0,b)$，即短轴上端点（或下端点）。因 $b>0$，代入方程验证：$\frac{0^{2}}{a^{2}}+\frac{b^{2}}{b^{2}}=1$，成立。
			\[
				P(0,b).
			\]
			\par\textbf{理由：}
			该点坐标满足椭圆方程，且对称，便于计算距离。

		\item[\textbf{步骤三（利用椭圆定义）：}]
			由椭圆定义，点 $P$ 到两焦点距离之和等于 $2a$：
			\[
				|PF_{1}|+|PF_{2}|=2a.
			\]
			\par\textbf{理由：}
			椭圆定义：平面内到两定点距离之和为常数 $2a$ 的点的轨迹。

		\item[\textbf{步骤四（计算距离）：}]
			计算 $|PF_{1}|$ 和 $|PF_{2}|$。由距离公式：
			\[
				\begin{aligned}
					|PF_{1}| &= \sqrt{(0-(-c))^{2}+(b-0)^{2}} \\
					&= \sqrt{c^{2}+b^{2}},\\[4pt]
					|PF_{2}| &= \sqrt{(0-c)^{2}+(b-0)^{2}} \\
					&= \sqrt{c^{2}+b^{2}}.
			\end{aligned}
			\]
			\par\textbf{理由:}
			代入坐标,注意 $P(0,b)$ 与 $F_{1}(-c,0)$、$F_{2}(c,0)$ 的坐标.

		\item[\textbf{步骤五(代入并化简):}]
			代入定义式:
			\[
				\sqrt{c^{2}+b^{2}}+\sqrt{c^{2}+b^{2}}=2a\quad\Longrightarrow\quad 2\sqrt{c^{2}+b^{2}}=2a.
			\]
			两边除以 $2$ 并平方:
			\[
				\sqrt{c^{2}+b^{2}}=a\quad\Longrightarrow\quad c^{2}+b^{2}=a^{2}.
			\]
			\par\textbf{理由:}
			直接代数运算,平方去掉根号,得到关系.

	\end{description}

	\medskip
	\textbf{\textcolor{CProof}{结论回扣}}

	\textbf{因此,对椭圆标准方程,总有 $a^{2}=b^{2}+c^{2}$,该式沟通了椭圆的几何参数,是后续计算离心率、焦点坐标的基础.}
\end{proofbox}
